% Source: https://artofproblemsolving.com/community/c3h2290078_modular_arithmetic_troll
% Stable identity: aops-topic-2290078-q01
\textbf{The error.} The rule ``if a product is $0$, one factor is $0$'' holds modulo a prime, but not modulo the composite number $14$: for example $2\cdot7\equiv0\pmod{14}$ with neither factor $\equiv0$. Indeed $x=7$ gives $7\cdot4=28\equiv0\pmod{14}$, a solution the proposed answer misses.

\textbf{Correct solution.} $14\mid x(x-3)$ exactly when $2\mid x(x-3)$ and $7\mid x(x-3)$.
\begin{itemize}
\item $x$ and $x-3$ have opposite parity, so $x(x-3)$ is always even.
\item $7$ is prime, so $7\mid x(x-3)$ exactly when $x\equiv0$ or $x\equiv3\pmod7$.
\end{itemize}
Modulo $14$ these are $\boxed{x\equiv0,\ 3,\ 7,\ 10\pmod{14}}$.
